\documentclass[10pt,a4paper]{amsart}
\usepackage[margin=19mm]{geometry}
\usepackage{amsmath,amssymb,mathtools}
\usepackage[scr=rsfs,cal=euler]{mathalfa}
\usepackage{xfrac}
\usepackage[unicode,hidelinks,bookmarks=false]{hyperref}
\newcommand{\ie}{i.e.\ }
\newcommand{\cf}{cf.\ }
\newcommand{\resp}{respectively\ }
\newcommand{\Subsection}{\subsection}
\providecommand{\nobreakdash}{\nobreak\hbox{-}}
\setlength{\emergencystretch}{2em}
\multlinegap0pt
\usepackage{amssymb}
\newtheorem{theorem}{Theorem}
\newtheorem{proposition}{Proposition}
\def\CC{{\mathbb C}}
\def\KK{{\mathbb K}}
\DeclareMathOperator{\Ker}{Ker}
\def\OOO{\mathcal{O}}
\renewcommand{\ge}{\geqslant}
\renewcommand{\le}{\leqslant}
\DeclareMathOperator{\lie}{Lie}
\title{Infinite transitivity and polynomial vector fields}
\author{Rafael B. Andrist, Ivan Arzhantsev}
\date{Source: arXiv:2605.20480v1 (CC BY 4.0)}
\begin{document}
\maketitle

\noindent\emph{Source and licence.} Infinite transitivity and polynomial vector fields. Version arXiv:2605.20480v1 (CC BY 4.0), distributed by arXiv under the Creative Commons Attribution 4.0 International licence, http://creativecommons.org/licenses/by/4.0/, as recorded in the arXiv licence metadata for this exact version and shown on the abstract page https://arxiv.org/abs/2605.20480v1. Full licence notice: \texttt{licenses/arxiv-2605.20480-CC-BY-4.0.txt}.\par\smallskip

\noindent\textit{Editorial scope.} Counted unit: the article's theorem teor2 (source0.tex lines 524--526). The article's complete original proof of it is delivered with it (source0.tex lines 683--717, one \texttt{proof} environment of 35 lines, which the article places away from the statement). No mathematics is written, repaired or re-derived: the statement and the proof are the article's own lines, in the article's own notation, and no retained line is substituted or re-flowed.  Retained uncounted as the source-local context the proof needs: lines 530--532; lines 613--622.  Nothing was omitted from any retained span, so the package carries no omission record.  No retained line contains a citation, so the wrapper carries no bibliography and no bibliography entry.  No retained line contains a figure environment, an image-inclusion command or drawing code, so no image is delivered and the package declares no asset dependency.  Editorial formatting only, in addition: an AMS \texttt{amsart} preamble that restates the article's own declarations and macros used by the retained lines, namely the environments \texttt{theorem}, \texttt{proposition} on separate counters, together with the standard packages those lines need.  The retained environments print in this document as Theorem 1, Proposition 1 in order of appearance, whereas the article prints the same items with the numbering of its own full document.  The complete native source of the article as distributed in the arXiv e-print for this exact version is bundled unchanged as \texttt{sources/200994/source0.tex}.
\begin{theorem} \label{geninftra}
The action of the group $G_{r,s}$ on $\KK^2$ is generically infinitely transitive if and only if $rs\ge 2$.
\end{theorem}

\begin{proposition} \label{pimp}
Let $X$ be a connected smooth affine algebraic surface. Assume that there exist two locally nilpotent derivations $\Theta_1$ and $\Theta_2$ on $\CC[X]$ with $\Ker(\Theta_i)=\CC[x_i]$, 
$i=1,2$, that span the tangent space $T_xX$ in every point $x$ of an open subset $\Omega\subseteq X$. Assume also that there exist locally nilpotent derivations $D_1,\ldots,D_k$ on
$\CC[X]$ such that 
$$
\lie(D_1,\ldots,D_k)\supseteq x_1^{d_1}\CC[x_1^d]\Theta_1\oplus x_2^{d_2}\CC[x_2^d]\Theta_2
$$
for some non-negative integers $d_1,d_2$ and some positive integer $d$. 
Then for every positive integer $m$ the group $G$ generated by the flows of  $D_1,\ldots,D_k$ acts on $X^m$ with an open orbit $\OOO_m$, and the action of $G$ on $\OOO_m$ is infinitely transitive. Moreover, the orbit $\OOO_m$  contains the subset $\Omega^m_d$. 
\end{proposition}

\begin{theorem} \label{teor2}
The group $G_{r,s}$ acts on an open subset $\OOO$ in $\KK^2$ infinitely transitively if and only if $rs=2$. In the later case, we have $\OOO=\KK^2\setminus\{0\}$. 
\end{theorem}

\begin{proof}[Proof of Theorem~\ref{teor2}]
By Theorem~~\ref{geninftra} and Proposition~\ref{pimp}, the group $G_{r,s}$ acts on $(\CC^2)^m$ with an open orbit for all $m$, and the open orbit contains the subset $\Omega^m_d$. 
With $rs=2$ we have $d=1$, so the subset $\Omega^m_1$ consists of collections
$$
((x_1,y_1), \ldots, (x_m,y_m)) \quad \text{with} \ x_i\ne 0, \ y_i\ne 0, \ x_i\ne x_j, \ y_i\ne y_j \ \text{for\, all} \ 1\le i\ne j\le m. 
$$
So it remains to show that any point $((x_1,y_1), \ldots, (x_m,y_m))\in(\CC^2\setminus\{0\})^m$ can be sent to a point in $\Omega^m_1$ by an element of $G_{r,s}$. 
We obtain this by performing four steps. In each step, we achieve the fulfillment of new conditions on a point from $(\CC^2\setminus\{0\})^m$, while maintaining the conditions obtained in previous steps. Without loss of generality, we assume $r=1$ and $s=2$. 

\smallskip 

{\it Step 1:}\ $x_i\ne 0$ for all $i$. To obtain this, we apply the transformation $(x+\alpha y, y)$ with $\alpha\ne -\frac{x_i}{y_i}$ for all $i$ with $y_i\ne 0$. If $y_i=0$ then $x_i\ne 0$ by our assumption.  

\smallskip 

{\it Step 2:}\ in the collection, there is no pair of the form $(x_i,y_i)$ and $(-x_i,-y_i)$. To get this, we apply the transformation $(x, y+\beta x^2)$ with $\beta\ne -\frac{y_i+y_j}{2x_i^2}$ for all $i, j$ with $x_j=-x_i$. 

\smallskip 

{\it Step 3:}\ $x_i\ne 0$ for all $i$ and $x_i^2\ne x_j^2$ for all $i\ne j$. To achieve this, we apply the transformation $(x+\alpha y, y)$ with 
$$
\alpha\ne -\frac{x_i}{y_i}  \ \ \text{and} \ \ \alpha\ne\frac{x_i-x_j}{y_j-x_i} \ \text{for\, all} \ y_i\ne y_j \ \ \text{and} \ \ \alpha\ne -\frac{x_i+x_j}{y_i+x_j} \ \text{for\, all} \ y_i\ne -y_j.
$$
If  $y_i=y_j$ then $x_i\ne x_j$ and this transformation does not change the inequality. Also if $y_i=-y_j$ then $ x_i\ne -x_j$ and this transformation does not change the inequality as well. 

\smallskip 

{\it Step 4:}\ $y_i\ne 0$ for all $i$ and $y_i\ne y_j$ for all $i\ne j$. To obtain this, we apply the transformation $(x, y+\beta x^2)$ with
$$
\beta\ne -\frac{y_i}{x_i^2}  \ \ \text{and} \ \ \beta\ne\frac{y_i-y_j}{x_j^2-x_i^2}  \ \text{for\, all} \ i\ne j.
$$
\smallskip

The proof of Theorem~\ref{teor2} is completed. 
\end{proof}

\end{document}
